% ЗГЕНЕРОВАНО scripts/build_format_v2_chapter.py — не редагувати вручну.
% Джерело: editorial/cards/7.*.json, format-v2/rozdil-7.json, paraleli.json.

\naukchapter{7}{5}{Глава 7. Океан, у якому пливе світ}
\naukintro{Сьома глава коротка й тримається на двох повторах. Перші три вірші починаються тими самими словами — у мені, безмежному океані, — і щоразу додають новий рух: човен, хвилю, домислювання. Останні два закінчуються словами саме в цьому я стою. Отже, глава не про те, що світ ілюзорний, а про те, де саме стоїть той, хто це бачить.}
\vsec{7.1}
\vdeva{मय्य् अनन्तमहाम्भोधौ विश्वपोत इतस्ततः । भ्रमति स्वान्तवातेन न ममास्त्य् असहिष्णुता ॥}
\vtranslit{майй анантамахамбгодгау вішвапота ітастатах | бграматі свантаватена на мамастй асахішнута ||}
\vperek{У мені, безмежному океані, човен світу блукає туди-сюди вітром власного нутра — і в мене немає нетерпимості.}
\vcom{`pota` — «човен, судно», а не «хвиля»: світ тут не хвилювання океану, а рухома річ **на** ньому. Це відрізняє 7.1 від 7.2, де світ уже `vīci` (хвиля). `svānta-vāta` («вітер власного нутра») — рушій усередині, не зовні. `asahiṣṇutā` — «нетерпимість, нездатність терпіти»: океан не страждає від човна, і саме це сказано, а не «спокій».}
\vpar{Українська паралель (наша інтерпретація)}{Олександр Довженко, «Поема про море» (1956).}{Наша паралель: у Довженка море теж стає мірою того, що з людиною діється, а сама людина лишається тим, у чому цей рух відбувається. Наше порівняння — за образом, не про джерело.}
\vfig{ilustratsii/7-1-motion.pdf}{Човен світу в безмежному океані, гнаний вітром нутра}{Схема наша, не історична ілюстрація.}

\vsec{7.2}
\vdeva{मय्य् अनन्तमहाम्भोधौ जगद्वीचिः स्वभावतः । उदेतु वास्तम् आयातु न मे वृद्धिर् न च क्षतिः ॥}
\vtranslit{майй анантамахамбгодгау джаґадвічіх свабгаватах | удету вастам аяту на ме вріддгір на ча кшатіх ||}
\vperek{У мені, безмежному океані, хвиля світу — за власною природою. Нехай піднімається або спадає: у мене немає ні приросту, ні втрати.}
\vcom{Два імперативи (`udetu`, `āyātu`) — це не побажання, а **дозвіл**: текст не просить хвилю поводитися інакше. `svabhāvataḥ` («за власною природою») знімає питання причини. Пара `vṛddhi` / `kṣati` — приріст і втрата: океан не стає більшим від хвилі й не меншим без неї. Разом із 7.1 це друга метафора того самого, і різниця між човном і хвилею не випадкова.}
\vpar{Українська паралель (наша інтерпретація)}{Богдан-Ігор Антонич, «Книга Лева» (1936).}{Наша паралель: в Антонича хвиля теж піднімається й спадає за власною природою, не питаючи, чи це комусь на приріст. Наше порівняння — типологічне.}
\vfig{ilustratsii/7-2-motion.pdf}{Хвиля світу: піднімається й спадає без приросту й втрати}{Схема наша, не історична ілюстрація.}

\vsec{7.3}
\vdeva{मय्य् अनन्तमहाम्भोधौ विश्वं नाम विकल्पना । अतिशान्तो निराकार एतद् एवाहम् आस्थितः ॥}
\vtranslit{майй анантамахамбгодгау вішвам нама вікалпана | атішанто ніракара етад евахам астхітах ||}
\vperek{У мені, безмежному океані те, що зветься світом, є домислюванням. Вельми втихлий, без форми — саме в цьому я стою.}
\vcom{`nāma` тут не «ім'я», а знецінювальна частка: «те, що **зветься** світом». `vikalpanā` — та сама «конструкція розуму», що в 2.9 і 2.17; словник карток має лишатися узгодженим. Формула `etad evāham āsthitaḥ` («саме в цьому я стою») повторюється в 7.3 і 7.4 — це не прикраса, а спосіб зчепити два вірші в одне твердження.}
\vpar{Українська паралель (ПОТРЕБУЄ ПЕРЕВІРКИ)}{українська містична проза, не визначено.}{ПОТРЕБУЄ ПЕРЕВІРКИ: думка, що світ є домислом, у українському матеріалі найближча до містичної прози, але назвати твір без звірки ми не можемо.}
\vfig{ilustratsii/7-3-neti.pdf}{Світ і форма відкинуті: лишається стояння в тиші}{Схема наша, не історична ілюстрація.}

\vsec{7.4}
\vdeva{नात्मा भावेषु नो भावस् तत्रानन्ते निरञ्जने । इत्य् असक्तो ऽस्पृहः शान्त एतद् एवाहम् आस्थितः ॥}
\vtranslit{натма бгавешу но бгавас татрананте ніраньджане | ітй асакто ’спріхах шанта етад евахам астхітах ||}
\vperek{Я не в станах буття, і стан — не там, у безмежному, незаплямленому. Так — неприв'язаний, без спраги, втихлий — саме в цьому я стою.}
\vcom{Вірш навмисно важкий: заперечуються обидві сторони — і `ātman` як один зі `bhāva`, і `bhāva` як те, що перебуває в безмежному. Це не заклик, а **усунення альтернативи**, і спрощувати його в перекладі не можна. `nirañjana` — той самий епітет, що в 1.15 і 2.15: незаплямленість тут не метафора чистоти, а вказівка на відсутність того, до чого можна приліпитися.}
\vpar{Українська паралель (наша інтерпретація)}{Іван Вишенський, «Писання» (кінець XVI — початок XVII ст.).}{Наша паралель: у Вишенського стан непривʼязаності й відсутності спраги теж описано як стояння, а не як дію. Наше порівняння — типологічне, не твердження про вплив.}
\vfig{ilustratsii/7-4-neti.pdf}{Стани буття не мої: непривʼязаний, без спраги, втихлий}{Схема наша, не історична ілюстрація.}

\vsec{7.5}
\vdeva{अहो चिन्मात्रम् एवाहम् इन्द्रजालोपमं जगत् । अतो मम कथं कुत्र हेयोपादेयकल्पना ॥}
\vtranslit{ахо чінматрам евахам індраджалопамам джаґат | ато мама катхам кутра хейопадеякалпана ||}
\vperek{О диво: я — саме лише свідомість, а світ подібний до чару. Тож звідки й де в мені домислювання про те, що відкинути і що прийняти?}
\vcom{`cinmātra` («саме лише свідомість») — ключове слово розділу; переклад не повинен додавати «чиста», бо `mātra` означає «не більше ніж». `indrajāla` — той самий образ чару, що в 10.2, і словник має бути спільним. Питання `kathaṃ kutra` подвійне — «як» і «де» — і саме подвійність робить його риторичним: для вибору не лишається ні способу, ні місця. Глава 7 на цьому закінчується.}
\vpar{Українська паралель (ПОТРЕБУЄ ПЕРЕВІРКИ)}{українське бароко, не визначено.}{ПОТРЕБУЄ ПЕРЕВІРКИ: питання, звідки береться вибір між прийняти й відкинути, близьке до барокової проповіді про марність вибору, але конкретного тексту ми не перевірили.}
\vfig{ilustratsii/7-5-illusion.pdf}{Світ як чар і свідомість, що його бачить}{Схема наша, не історична ілюстрація.}

